\textcite{Beck2004XP} define Agile as a set of values and principles that emphasize iterative development, customer collaboration, and responsiveness to change. Agile development approaches typically include iterative and incremental development, self-organizing teams, and continuous feedback and adaptation.

Agile processes focus on the rapid and early delivery of working software and are based on development models that support incremental and iterative progress \parencite{keith2010agile}.

A user story is a description of a feature expressed from the perspective of the user and its intended value. User stories serve as the core unit of work in Agile development. Multiple user stories are maintained in a backlog, from which a subset is selected for implementation during a time-boxed iteration, typically lasting two to four weeks, known as a sprint. A collection of related user stories that collectively deliver a larger functional capability is referred to as an epic. Completed sets of functionality, such as alpha or beta versions, are commonly referred to as milestones.

User stories must be estimated, which requires sufficient understanding of both the system being built and the implementation approach. If the scope of a story is too large or insufficiently understood, accurate estimation becomes difficult. Such cases are often addressed through spikes, which are time-boxed research tasks used to reduce uncertainty \parencite{keith2010agile}.

Both DBR and Agile share a common iterative philosophy, in which development proceeds through cycles of analysis, implementation, evaluation, and refinement, enabling continuous learning and adaptation throughout the process.